\section{Finite memory certificates without physical decoder assumptions}
\label{sec:r77-unrestricted-approximation}

The physical-dictionary lower bound uses the geometry of density matrices
at individual hidden labels.  An arbitrary classical simulator need not
have that geometry.  We now give a different certificate, valid for
arbitrary classical decoder tables.  Its conclusion is quantitatively
weaker, but applies to the original algebraic qubit instrument without
an unproved Diophantine hypothesis.  More generally it applies whenever
experimentally accessible preparations have separated exposing tests.

Nonnegative factorizations, support antichains and their relation to
memory and extension complexity are classical.  In particular the exact
zero-diagonal support bound is already in the literature on nonnegative
rank; see \cite[Section 4.1 and Corollary 4.13]{GPT13} and
\cite[Section 3.1]{VandaeleGillisGlineur17}.  The argument below supplies
an explicit error margin and derives its hypotheses from joint stopped
records.  It does not treat an approximate zero as an exact zero, or
assume that a rare conditional distribution inherits the unconditional
approximation error.

\subsection{A robust support certificate}

For $J\ge2$ define the antichain threshold
\begin{equation}\label{eq:r77-antichain-threshold}
 a(J)=\min\left\{m\ge1:
                   \binom{m}{\lfloor m/2\rfloor}\ge J\right\}.
\end{equation}
In particular $\lceil\log_2J\rceil\le a(J)\le J$.

\begin{lemma}[Robust antichain factorization certificate]
\label{lem:r77-robust-antichain}
Let $\lambda_1,\ldots,\lambda_J$ be probability rows on $M$ labels,
and let $d_1,\ldots,d_J\in[0,1]^M$ be arbitrary columns.  If
\begin{equation}\label{eq:r77-antichain-hypotheses}
 \lambda_n d_n\le\epsilon,
 \qquad \lambda_n d_m\ge b\quad(n\ne m),
 \qquad 0<b\le1,
 \qquad \epsilon<\frac{b^2}{2M^2},
\end{equation}
then $J\le\binom{M}{\lfloor M/2\rfloor}$.
The conclusion uses no linearity, noncontextuality or quantum
representation of the decoder columns.
\end{lemma}
\begin{proof}
Put $S_n=\{i:\lambda_{n,i}\ge b/(2M)\}$.  For $n\ne m$, the
nonnegative sum $\lambda_nd_m\ge b$ has a summand
$\lambda_{n,i}d_m(i)\ge b/M$.  Since both factors are at most one,
\[
 \lambda_{n,i}\ge b/M,\qquad d_m(i)\ge b/M.
\]
Thus $i\in S_n$.  On the other hand
\[
 \lambda_{m,i}\le\epsilon/d_m(i)
                \le M\epsilon/b<b/(2M),
\]
so $i\notin S_m$.  Every ordered pair of distinct sets has this
witness, and the $S_n$ form an antichain in the subsets of $[M]$.

For completeness, choose a uniformly random permutation of $[M]$.
For a fixed set $S$ of size $k$, the probability that its elements
are the first $k$ elements of this permutation, in any order, is
$1/\binom{M}{k}$.  The corresponding events for members of an
antichain are disjoint, since the initial sets of one permutation
form a chain.  Hence
\[
 \frac{J}{\binom{M}{\lfloor M/2\rfloor}}
 \le\sum_{n=1}^J\binom{M}{|S_n|}^{-1}\le1,
\]
as required.
\end{proof}

The combinatorial conclusion is sharp at the level of these hypotheses.
For $M\ge2$, let $k=\lfloor M/2\rfloor$, index preparations and
tests by all $k$-element subsets $S$ of $[M]$, and put
\[
 \lambda_S(i)=\frac{\mathbf1_{\{i\in S\}}}{k},
 \qquad d_T(i)=\mathbf1_{\{i\notin T\}}.
\]
Then $\lambda_Sd_S=0$ and $\lambda_Sd_T\ge1/k$ for $S\ne T$, with
exactly $\binom{M}{k}$ preparations.  This construction concerns the
abstract preparation--test table.  It does not construct a small
controller for the qubit's chronological command--report dynamics.

\subsection{An operational certificate from joint records}

Consider any causal experiment with $J$ preparation procedures.  The
$n$th procedure either declares failure or ends at an observable herald
$H_n$.  Only after that herald is a binary terminal test $m\in[J]$
requested.  The choice of $m$, including a policy's future test code,
is not an input to the simulator before the herald.  Its terminal decoder
is one fixed table for each requested
test, independent of the earlier preparation procedure except through
the simulator's current counted label.  These are the usual causal
finite-memory conventions; allowing a decoder to inspect an uncounted
copy of the entire past would define a different resource.

Let $p_n$ be the target probability of $H_n$, and let $f_{nm}$ be the
target rejection probability conditional on $H_n$ and test $m$.
Suppose that every one of these preparation--test experiments has
target expected cost at most $B$, and that
\begin{equation}\label{eq:r77-operational-exposure}
 p_n\ge p_*>0,\qquad f_{nn}=0,\qquad
 f_{nm}\ge s>0\quad(n\ne m).
\end{equation}
The terminal tests and preparation procedures may be classical or
quantum; only their observable probabilities enter the next theorem.

\begin{theorem}[Memory lower bound from approximately exposing tests]
\label{thm:r77-operational-antichain}
Suppose \eqref{eq:r77-operational-exposure} holds, with $J\ge2$ and
$0<s\le1$.  An $M$-label classical simulator that approximates the
complete stopped record of every specified experiment to total-variation
error at most $\delta$ satisfies
\begin{equation}\label{eq:r77-budgeted-antichain}
 \delta\le\frac{p_*s^2}{32M^2}
 \quad\Longrightarrow\quad
 J\le\binom{M}{\lfloor M/2\rfloor}.
\end{equation}
Consequently
\begin{equation}\label{eq:r77-operational-lower}
 \delta\le\frac{p_*s^2}{32J^2}
 \quad\Longrightarrow\quad M\ge a(J).
\end{equation}
Both assertions allow freely mixed or counted pure starts and arbitrary
stochastic terminal decoder tables.  They require only the finite
collection of $J^2$ displayed preparation--test experiments.
\end{theorem}
\begin{proof}
For the $n$th procedure let $q_n$ be the simulator's herald probability,
and let $\lambda_n$ be its conditional distribution on its counted
labels at that herald.  It does not depend on the later requested $m$.
Write $d_m(i)\in[0,1]$ for the rejection probability of test $m$ at
label $i$.  Total variation controls both the herald event and the
joint herald-and-rejection event, so
\begin{equation}\label{eq:r77-joint-error}
 |q_n-p_n|\le\delta,\qquad
 |q_n\lambda_nd_m-p_nf_{nm}|\le\delta.
\end{equation}
Under the hypothesis in \eqref{eq:r77-budgeted-antichain},
$\delta/p_n\le s^2/(32M^2)\le1/32$, so $q_n>0$ and
\[
 \lambda_nd_n\le\frac{\delta}{p_n-\delta}
                 \le\frac{2\delta}{p_*}
                 \le\frac{s^2}{16M^2}.
\]
For $m\ne n$ put $t=\delta/p_n\le s^2/(32M^2)\le s/4$.
Equation~\eqref{eq:r77-joint-error} gives
\[
 \lambda_nd_m\ge\frac{p_ns-\delta}{p_n+\delta}
                =\frac{s-t}{1+t}\ge\frac{s}{2}.
\]
Apply Lemma~\ref{lem:r77-robust-antichain} with $b=s/2$ and
$\epsilon=s^2/(16M^2)<b^2/(2M^2)$.  This proves
\eqref{eq:r77-budgeted-antichain}.  If $M\ge J$, then $M\ge a(J)$
already.  Otherwise the hypothesis in \eqref{eq:r77-operational-lower}
implies the one in \eqref{eq:r77-budgeted-antichain}, and gives the
same conclusion.  No conditional error estimate was presumed before
\eqref{eq:r77-joint-error}; normalization was performed only after
the herald mass was controlled.
\end{proof}

This theorem keeps the cost budget in the target law.  It makes no
assertion about the simulator's expected duration.  The hidden label
can have any classical interpretation, and the decoder for one binary
POVM can be unrelated to decoders for other POVMs.

\subsection{An unrestricted lower bound for the original algebraic instrument}

For the undamped qubit instrument
\eqref{eq:r75-noisy-instrument}, fix one commanded Bloch angle
$\theta$ and put $z=e^{i\theta}$ and $\omega=\theta/(2\pi)$.  Let
$N_{\rm cl}^{\rm m}(\delta;B)$ and
$N_{\rm cl}^{\rm p}(\delta;B)$ be the least label counts approximating
every common visible stopping and terminal-POVM policy of target expected
calls at most $B$, with arbitrary classical terminal decoder tables.
Transitions are fixed independently of later terminal requests.  Thus
these classes include the physical dictionaries of
Section~\ref{sec:r76-physical-approximation}, but impose no Born
representation on their individual labels.

Fix $B>0$ and write $\zeta=\min\{1,7B/8\}$.  For $J\ge2$ and
a phase which is not a root of unity put
\begin{equation}\label{eq:r77-phase-margin}
 s_J=\min_{1\le q<J}\frac{1-\operatorname{Re}(z^q)}{2}>0.
\end{equation}

\begin{theorem}[Unrestricted finite certificates and arithmetic rates]
\label{thm:r77-unrestricted-qubit}
For either start convention, the following statements hold.

\smallskip\noindent
\textup{(i)} For every non-root phase and every integer $J\ge2$,
\begin{equation}\label{eq:r77-phase-memory}
 0<\delta\le\frac{\zeta\,8^{-J}s_J^2}{32J^2}
 \quad\Longrightarrow\quad N_{\rm cl}(\delta;B)\ge a(J).
\end{equation}
This lower bound uses $J$ terminal projectors and finitely many
preparation prefixes, with no assumption on individual decoder rows.

\smallskip\noindent
\textup{(ii)} At the original algebraic phase
$z=(3+4i)/5$, the explicit sufficient condition is
\begin{equation}\label{eq:r77-rational-phase-memory}
 0<\delta\le\frac{25\zeta}{128\cdot200^J J^2}
 \quad\Longrightarrow\quad N_{\rm cl}(\delta;B)\ge a(J).
\end{equation}
No badly-approximable assertion about $\theta/(2\pi)$ is required.

\smallskip\noindent
\textup{(iii)} If $\theta=2\pi\omega$ and
$\|q\omega\|_{\mathbb R/\mathbb Z}\ge c/q$ for every $q\ge1$,
then another explicit sufficient condition is
\begin{equation}\label{eq:r77-badly-approx-memory}
 0<\delta\le\frac{\zeta c^4}{2\cdot8^J J^6}
 \quad\Longrightarrow\quad N_{\rm cl}(\delta;B)\ge a(J).
\end{equation}

\smallskip\noindent
In case \textup{(ii)}, and also in case \textup{(iii)},
\begin{equation}\label{eq:r77-unrestricted-growth}
 \liminf_{\delta\downarrow0}
 \frac{N_{\rm cl}^{\rm m}(\delta;B)}
      {\log_2\log_8(1/\delta)}\ge1,
 \qquad
 \liminf_{\delta\downarrow0}
 \frac{N_{\rm cl}^{\rm p}(\delta;B)}
      {\log_2\log_8(1/\delta)}\ge1.
\end{equation}
The inherited physical constructions give, at the same fixed phase,
$N_{\rm cl}^{\rm m}\le K_\omega+L-1$ and
$N_{\rm cl}^{\rm p}\le K_\omega+L$, where $L$ is
\eqref{eq:r76-qubit-depth}; at $z=(3+4i)/5$ one may take $K_\omega=7$.
The unrestricted lower and upper therefore leave a gap between
logarithmic-in-logarithm and logarithmic growth.  A matching logarithmic
lower is asserted only for the physical-dictionary class under the
hypotheses of Theorem~\ref{thm:r76-physical-logarithm}.
\end{theorem}
\begin{proof}
For $n=0,\ldots,J-1$, begin with probability $\zeta$ and otherwise
stop with no call.  If begun, follow the report prefix $1\,0^n$,
stopping immediately on a different report.  Success is the herald
$H_n$.  The target probability and conditional terminal state are
\[
 p_n=\zeta8^{-(n+1)},\qquad
 \rho_n=U_\theta^n\tau(U_\theta^*)^n.
\]
In the real Bloch plane $\rho_n$ has vector
$(\cos(n\theta),\sin(n\theta))$.  The target expected number of
calls, before the cost-free terminal test, is
\[
 \zeta\sum_{j=0}^{n}8^{-j}\le\zeta\,8/7\le B.
\]
After success request the binary terminal POVM
$(\rho_m,I-\rho_m)$, and call its second outcome rejection.  Its
target conditional rejection probability is
\[
 f_{nm}=1-\tr(\rho_n\rho_m)
       =\frac{1-\cos((n-m)\theta)}2.
\]
Thus $f_{nn}=0$, $f_{nm}\ge s_J$ for $n\ne m$, and
$p_n\ge p_*:=\zeta8^{-J}$.  These are the hypotheses of
Theorem~\ref{thm:r77-operational-antichain}, which gives (i).
The simulator's terminal projector request is made only after the
prefix.  Causality is precisely what gives one common label
distribution $\lambda_n$ for all its $J$ possible terminal tests.

For (ii), write
\[
 z^q=\frac{A_q+iB_q}{5^q},\qquad A_q,B_q\in\mathbb Z.
\]
The phase is not a root of unity: otherwise $z+z^{-1}=6/5$ would
be both a rational number and an algebraic integer, hence an integer.
Consequently $\operatorname{Re}z^q<1$ for every $q\ge1$, and
$5^q-A_q$ is a positive integer.  Therefore
\[
 \frac{1-\operatorname{Re}z^q}{2}
 =\frac{5^q-A_q}{2\cdot5^q}\ge\frac1{2\cdot5^q},
 \qquad s_J\ge\frac1{2\cdot5^{J-1}}.
\]
Substitution in (i) gives
\[
 \frac{\zeta8^{-J}}{32J^2}\frac1{4\cdot25^{J-1}}
 =\frac{25\zeta}{128\cdot200^J J^2},
\]
which proves the claimed rational certificate.  The phase here is
the Bloch rotation, as in \eqref{eq:r75-noisy-instrument}; its unitary
uses the half-angle.  No angle is squared in this calculation.

For (iii), $\sin(\pi t)\ge2t$ for $0\le t\le1/2$ gives,
for $1\le q<J$,
\[
 \frac{1-\cos(2\pi q\omega)}2
 =\sin^2(\pi q\omega)
 \ge4\|q\omega\|_{\mathbb R/\mathbb Z}^2
 \ge\frac{4c^2}{q^2}\ge\frac{4c^2}{J^2}.
\]
Part (i) then yields \eqref{eq:r77-badly-approx-memory}.

For fixed $B$, the sufficient condition in (ii) holds with
$J=\log_{200}(1/\delta)-O_B(\log\log(1/\delta))$ chosen as an
integer slightly below this expression.  The corresponding choice in
(iii) is $J=\log_8(1/\delta)-O_{B,c}(\log\log(1/\delta))$.
Since $a(J)\ge\log_2J$, both imply
\eqref{eq:r77-unrestricted-growth}.  The asserted uppers are the
unchanged orbit-prefix constructions of
Theorem~\ref{thm:r76-physical-logarithm} and its final algebraic-phase
paragraph, which are also admissible in the unrestricted class.
\end{proof}

In fact the exact integer lower $a(J)$ retains a second-order term:
the central-binomial estimate
$\binom{m}{\lfloor m/2\rfloor}=\Theta(2^m/\sqrt m)$ gives
$a(J)=\log_2J+\tfrac12\log_2\log_2J+O(1)$.
This is a refinement of the support certificate, not an assertion of
the optimal asymptotic simulator size.  The lower uses terminal
projector requests, whereas the earlier exact infinite-memory
obstruction already used report words alone.  These are different
sets of operational probes.  No positive extra initializer charge is
deduced for the unrestricted approximation problem.

\begin{corollary}[Every non-root algebraic phase has the same lower order]
\label{cor:r77-algebraic-unrestricted}
Let $z$ be algebraic with $|z|=1$ and not a root of unity.  Let its
degree be $D$, choose an integer $A\ge1$ such that $Az$ is an
algebraic integer, and enumerate its conjugates as
$z=z_1,z_2,\ldots,z_D$.  Choose any supplied real bound
\[
 H\ge\max\left\{2,
             A^D\prod_{j=2}^D\max(1,|z_j|)\right\}.
\]
At that fixed phase, for either start convention,
\begin{equation}\label{eq:r77-algebraic-margin}
 0<\delta\le
 \frac{\zeta H^4}{32\cdot16^D\,(8H^4)^J J^2}
 \quad\Longrightarrow\quad N_{\rm cl}(\delta;B)\ge a(J).
\end{equation}
In particular \eqref{eq:r77-unrestricted-growth} holds.  No
badly-approximable premise is necessary for this lower order.
\end{corollary}
\begin{proof}
The number $A^q(z^q-1)=(Az)^q-A^q$ is a nonzero algebraic integer.
Its norm is a nonzero integer, so
\[
 1\le A^{qD}|z^q-1|\prod_{j=2}^D|z_j^q-1|
   \le 2^{D-1}H^q|z^q-1|.
\]
Because $|z|=1$,
\[
 \frac{1-\operatorname{Re}z^q}{2}
   =\frac{|1-z^q|^2}{4}\ge4^{-D}H^{-2q}.
\]
It follows that $s_J\ge4^{-D}H^{-2(J-1)}$.
Substitute this in \eqref{eq:r77-phase-memory} to obtain
\eqref{eq:r77-algebraic-margin}.  A choice
$J=\log_{8H^4}(1/\delta)-O_{B,H,D}(\log\log(1/\delta))$
then gives the claimed lower order.  This elementary algebraic norm
bound does not supply the stronger $s_J\asymp J^{-2}$ bound used for
the sharp physical-dictionary logarithmic law.
\end{proof}
